% !TEX root = ../main.tex

\section{Comparaison des deux méthodes}

La NMF et K-means ont été appliqués au même corpus et à la même représentation TF-IDF, mais ils n'en proposent pas la même lecture. La NMF fait apparaître des topics latents auxquels chaque segment est associé avec des poids variables. K-means répartit au contraire les segments en groupes lexicaux distincts. Leur comparaison permet d'identifier les régularités qui se retrouvent dans les deux organisations du corpus, ainsi que les différences liées au fonctionnement de chaque méthode.

\subsection{Deux représentations complémentaires}

La NMF offre une représentation graduelle des thèmes. Un même segment peut être associé à plusieurs topics, par exemple la famille, l'argent et l'espace domestique, avec une importance différente pour chacun. Cette souplesse correspond assez bien à la complexité des textes littéraires, dans lesquels plusieurs dimensions thématiques peuvent coexister.

K-means produit une organisation plus catégorielle : chaque segment est rattaché à un seul cluster, celui dont il est le plus proche dans l'espace vectoriel. Cette méthode fait ainsi apparaître des groupes lexicaux plus nettement séparés, mais rend moins visible la coexistence de plusieurs thèmes dans un même passage.

Les deux approches ne sont donc pas concurrentes. La NMF décrit la composition thématique des segments, tandis que K-means propose une partition générale du corpus. Leur rapprochement constitue une triangulation exploratoire, mais pas une validation indépendante, puisque les deux modèles utilisent les mêmes données et la même matrice TF-IDF.

\subsection{Convergences lexicales et thématiques}

La comparaison des termes caractéristiques fait apparaître plusieurs ensembles communs aux topics NMF et aux clusters K-means. Les deux méthodes repèrent notamment des motifs liés au travail et à la mine, à l'argent et à la finance, à la religion, à la famille, à la guerre, au commerce, aux espaces urbains et à la vie domestique.

Ces rapprochements montrent que certains champs lexicaux occupent une place structurante dans le corpus. Ils ne signifient toutefois pas que les résultats sont identiques. Un topic NMF peut réunir plusieurs dimensions que K-means répartit entre différents clusters. À l'inverse, un cluster peut rassembler des segments lexicalement proches dont l'interprétation thématique reste assez large.

La convergence entre les méthodes doit donc être appréciée à l'échelle des grands motifs plutôt qu'à travers une correspondance exacte entre chaque topic et chaque cluster. Comme les deux analyses reposent sur le même prétraitement, elle ne suffit pas à prouver l'existence de catégories thématiques objectives. Elle indique néanmoins que certains motifs résistent au changement de méthode de regroupement.

\subsection{Distribution des topics NMF par œuvre}

Pour étudier la répartition des topics dans les romans, chaque segment a été associé au topic possédant le poids NMF le plus élevé. Une table croisée entre les œuvres et les topics dominants a ensuite été construite et normalisée par ligne. Cette normalisation indique, pour chaque roman, la proportion de segments dominés par chacun des 29 topics.

\begin{figure}[H]
    \centering
    \includegraphics[width=\textwidth]{images/chapitre_1/resultats/proportion_segment_headmap_nmf.png}
    \caption{Proportion des segments par topic NMF dominant et par œuvre}
    \label{fig:distribution_topics_nmf_oeuvre}
\end{figure}

La figure~\ref{fig:distribution_topics_nmf_oeuvre} fait apparaître des associations attendues entre certaines œuvres et certains thèmes. Le monde de la mine et du travail ouvrier occupe ainsi une place importante dans \textit{Germinal}, tandis que la finance ressort davantage dans \textit{L'Argent}. Le commerce est particulièrement visible dans \textit{Au Bonheur des dames}, et la guerre dans \textit{La Débâcle}.

D'autres topics sont répartis entre plusieurs romans et correspondent à des motifs plus transversaux, comme la famille, les relations sociales, la religion ou l'espace domestique. La normalisation par œuvre évite que les romans les plus longs soient automatiquement surreprésentés et facilite ainsi la comparaison de leurs profils thématiques.

Cette visualisation repose néanmoins sur le topic dominant de chaque segment. Elle simplifie donc la représentation fournie par la NMF, puisqu'elle ne montre pas les topics secondaires également présents dans les passages.

\subsection{Distribution des clusters K-means par œuvre}

Une seconde heatmap présente la répartition des clusters K-means dans chaque roman. Comme pour la NMF, les valeurs sont exprimées sous forme de proportions afin de neutraliser autant que possible les différences de longueur entre les œuvres.

\begin{figure}[H]
    \centering
    \includegraphics[width=\textwidth]{images/chapitre_1/resultats/headmap_kmeans.png}
    \caption{Distribution des clusters K-means par œuvre}
    \label{fig:distribution_clusters_kmeans_oeuvre}
\end{figure}

La figure~\ref{fig:distribution_clusters_kmeans_oeuvre} fait apparaître plusieurs associations proches de celles obtenues par la NMF. \textit{Germinal} se distingue par le cluster consacré à la mine, aux ouvriers et à la grève. \textit{L'Argent} est fortement associé à la finance et aux affaires, tandis que \textit{La Débâcle} se caractérise par le vocabulaire des soldats, de l'armée et de la guerre. Le cluster lié au magasin, à la vente et aux clients est particulièrement présent dans \textit{Au Bonheur des dames}.

Ces rapprochements confirment que plusieurs œuvres possèdent un profil lexical fortement lié à leur univers narratif. La comparaison révèle également des différences : K-means impose un seul cluster à chaque segment, alors que la NMF peut y détecter plusieurs composantes thématiques. La distribution obtenue par K-means paraît donc généralement plus tranchée.

\subsection{Regroupement en grandes familles}

Afin de rendre les 29 clusters plus faciles à interpréter, une classification hiérarchique a été appliquée à leurs centres. Elle permet de les réunir en cinq familles thématiques plus générales.

\begin{figure}[H]
    \centering
    \includegraphics[width=0.85\textwidth]{images/chapitre_1/resultats/grande_famille_kmeans.png}
    \caption{Distribution des grandes familles de clusters K-means par œuvre}
    \label{fig:distribution_familles_kmeans_oeuvre}
\end{figure}

La figure~\ref{fig:distribution_familles_kmeans_oeuvre} synthétise la distribution de ces familles dans les romans. Elle distingue les motifs liés à l'argent et à la circulation financière, aux conflits politiques et militaires, aux institutions et à l'encadrement social, à la religion et à l'autorité cléricale, ainsi qu'à la vie quotidienne et aux expériences individuelles.

Cette représentation simplifie la lecture générale du corpus. La famille consacrée à la vie quotidienne est présente dans la majorité des œuvres, tandis que les autres ensembles apparaissent de manière plus concentrée. Les motifs financiers sont notamment visibles dans \textit{L'Argent}, les conflits militaires dans \textit{La Débâcle}, et la religion dans les romans où l'autorité ecclésiastique occupe une place centrale.

\subsection{Apport de la double approche}

La comparaison entre NMF et K-means permet finalement d'articuler deux niveaux d'analyse. La NMF rend compte du caractère souvent multiple des thèmes présents dans un segment, tandis que K-means propose une organisation plus nette des proximités lexicales. Leur rapprochement aide ainsi à distinguer les motifs récurrents des regroupements davantage liés au fonctionnement d'un modèle particulier.

Lorsque les résultats convergent, ils renforcent l'hypothèse qu'un motif occupe une place importante dans le corpus zolien. Lorsqu'ils divergent, l'examen des segments concernés permet de mieux comprendre les limites de chaque représentation. Ces régularités doivent cependant rester considérées comme des caractéristiques internes à l'œuvre de Zola. Leur éventuelle association au naturalisme devra être évaluée à partir du corpus comparatif et de la classification automatique.